\documentclass[12pt,a4paper]{article}
\usepackage{fontspec}
\setmainfont{Times New Roman}
\setmonofont{Courier New}
\usepackage[russian]{babel}
\usepackage{amsmath,geometry,booktabs,longtable,hyperref}
\geometry{margin=22mm}
\emergencystretch=3em
\title{Руководство по эксперименту и обработке видеозаписей шариков\\Версия 2.4}
\author{И. Л. Литвак}
\date{}
\begin{document}
\maketitle
Я подготовил эту редакцию для работы с текущим программным контуром \path{Mashine_Vision_Balls}. Она заменяет рабочую инструкцию, направленную 2 апреля 2026 года, в части порядка обработки, калибровки и состава результатов. Исходная инструкция сохраняется отдельно. Руководство устанавливает порядок будущих измерений; его наличие не подтверждает проведение описанных опытов.

\textbf{Программное обеспечение.} Реализованы пакетный анализ нескольких объектов, сопровождение, скалярная калибровка и экспорт таблиц, графиков и изображений. Опробовано подключение камеры. Метрологические характеристики определяются при экспериментальной проверке стенда.

Основной вход текущей Python-программы --- сохранённое видео. Запись через промышленную камеру выполняется отдельно средствами её SDK или программой регистрации. Наличие другого приложения в репозитории не означает, что его функции захвата встроены в описываемый пакетный анализ. Интерактивный режим текущего Python-контура ограничен прежним режимом одного шарика.

\section{Комплектация и подготовка оптической системы}
\begin{longtable}{p{61mm}p{70mm}}
\toprule Позиция & Назначение\\\midrule\endhead
Daheng MER2-231-41GM & Монохромная камера GigE; Sony IMX249, глобальный затвор, $1920\times1200$ пикселей, паспортно до 41 кадра/с\\
\path{VT-LEM1214CB-H1} & Объектив Contrastech, 12 мм, F1.4, C-mount, формат 1 дюйм\\
\path{VT-C2CS} & Кольца 0.5, 1, 5, 10 и 20 мм для подбора рабочей геометрии\\
\path{VT-LT2-HR12076-90W} & Кольцевой белый осветитель с диффузором, 24 В\\
\path{VT-LT4-2420PWDC-2} & Двухканальный контроллер освещения, 24 В, 20 Вт\\
\path{HR25-7TP-8S}, крепление NS & Кабель I/O камеры и крепление с резьбой 1/4 дюйма\\
\path{LT2-CB-1}, \path{GST25E12-P1J} & Кабель осветителя и адаптер MEAN WELL 12 В, 25 Вт; назначение адаптера в фактической схеме фиксируется отдельно\\
Raylab RL-GT3, RL-GT2 & Механические держатели\\\bottomrule
\end{longtable}
Состав приведён по закупочным документам, а не по акту осмотра собранного стенда. В старой инструкции указано 12.5 мм для объектива; в новой я использую 12 мм согласно счёту № 15022. До опыта необходимо сверить маркировку установленного объектива и фактическую комплектацию. Для кабеля подсветки следует записать фактическую длину: поставщик сообщал о временном трёхметровом кабеле вместо заказанного метрового.

Камера уже формирует монохромные изображения. Преобразование цветного сигнала в серый в её аппаратном тракте не требуется. При чтении видео OpenCV может вернуть трёхканальное представление; программа переводит его в один яркостный канал для единого входа алгоритма.

Перед подачей питания следует сверить маркировку устройств, напряжения и распиновку с документацией изготовителя. Осветитель на 24 В нельзя подключать к 12-вольтовому адаптеру только на основании наличия этого адаптера в закупке. Схему питания, подключения и местные правила работы с криогенными средами необходимо согласовать с ответственным за стенд.

\subsection{Монтаж и настройки}
Камеру и ячейку закрепляют жёстко. Подбирают рабочее расстояние, кольца, диафрагму и фокус так, чтобы одновременно различались контур шарика и эталон в рабочей плоскости. Фокусировку проверяют по центру и краям области измерения. Освещение подбирают без насыщения пикселей на измеряемой границе. Экспозицию и усиление фиксируют; автоматические изменения настроек при метрологической серии отключают, если они меняют видимость края.

Для каждого опыта записывают разрешение и область считывания, фактическую частоту и временные метки, экспозицию, усиление, диафрагму, режим освещения, рабочее расстояние, положение эталона и уровень жидкости. Сохраняют фотографию монтажа и настройки камеры. Паспортные 41 кадр/с не подставляют вместо измеренной частоты регистрации.

\section{Геометрия и эталонная серия}
Числа 55 мм, 38.13 мм, 0.85 мм и 6.3 мм в старой инструкции относятся к исходной геометрии: расстоянию шарик--чехол, расстоянию ячейка--чехол, толщине чехла и зазору соответственно. Они не являются протоколом новой серии. Значение 38.13 мм нельзя использовать как внутренний диаметр ячейки. Все необходимые размеры измеряют заново и сохраняют с неопределённостями.

Предпочтительный рабочий вариант --- эталон известного диаметра в плоскости шариков при тех же окне, жидкости, освещении и настройках камеры. Эталон должен иметь подтверждённый размер и чистый измеряемый контур. Изменение высоты шарика относительно эталона вызывает систематическое смещение масштаба. Скалярная калибровка текущей программы не исправляет дисторсию, преломление и перспективные искажения.

Порядок калибровки:
\begin{enumerate}
\item Снять отдельное видео одного эталона; проверить отсутствие движения камеры, насыщения и перекрытий.
\item Выполнить предварительный анализ. Полученные миллиметры при непроверенном начальном масштабе не использовать как измеренный диаметр.
\item Проверить аннотированные кадры и выбрать наблюдения только одного эталонного трека. CSV должен содержать столбец \texttt{ball\_diameter\_px}, минимум два положительных значения и один \texttt{ball\_id}. Сохранить выборку отдельно, не изменяя исходный экспорт.
\item Создать JSON масштаба по известному диаметру и его стандартной неопределённости.
\item Повторно обработать контрольное видео с полученным JSON. Для независимой проверки использовать другой эталон или независимую серию; восстановление диаметра калибровочного объекта не доказывает точность на других объектах.
\end{enumerate}
\begin{verbatim}
py -3 run_analysis.py --video data/videos/reference.mp4
  --output-dir output/reference_initial --frame-step 1

py -3 run_analysis.py --calibrate-csv reference_one_track.csv
  --calibration-diameter-mm 4.000
  --calibration-sigma-mm 0.005
  --calibration-width-px 1920
  --calibration-output calibration.json
\end{verbatim}
Команды показаны с переносами для чтения; при запуске каждую команду вводят одной строкой. Значения 4.000 мм, 0.005 мм и 1920 --- пример синтаксиса. Их заменяют действительными параметрами эталона и файла. Ширина должна соответствовать исходному видео, а не уменьшенной копии для детекции.

Масштаб и его модельная стандартная неопределённость определяются как
\[
S=\frac{d_{\mathrm{ref}}}{\overline{N}_{\mathrm{ref}}},\qquad
u(S)=S\sqrt{\left(\frac{u(d_{\mathrm{ref}})}{d_{\mathrm{ref}}}\right)^2+
\left(\frac{s(N_{\mathrm{ref}})}{\overline{N}_{\mathrm{ref}}}\right)^2}.
\]
Программа сохраняет наблюдаемый разброс кадров без деления на $\sqrt n$, поскольку последовательные кадры могут быть коррелированы. JSON применим только при неизменной оптической геометрии. Проверка ширины изображения программой не обнаруживает смену объектива, высоты или среды; за это отвечает протокол опыта.

\section{Запись экспериментальной серии}
Сначала проводят сухую контрольную серию на эталонах. Затем, при разрешённой работе с криогенной средой, повторяют контроль в рабочей оптической среде. До начала основной серии проверяют фокус, поле зрения, яркость и свободное место для записи.

Каждая серия получает уникальный идентификатор. В журнале фиксируют дату, оператора, материал шариков, способ получения, режим подачи, состояние поверхности, уровень жидкости, температуру при наличии измерения и все изменения настроек. Исходное видео сохраняют без повторного кодирования и обрезки. Производные фрагменты хранят отдельно со ссылкой на оригинал. После перемещения камеры, смены разрешения, оптики, фокуса или уровня, влияющего на масштаб, получают новую контрольную серию и проверяют применимость калибровки.

\section{Хранение и анализ через ClickHouse}
Основной запуск требует доступного сервера ClickHouse. Подключение реализовано через официальный HTTP-интерфейс \url{https://clickhouse.com/docs/interfaces/http}. Локальный сервер разворачивается из \path{infra/clickhouse.compose.yaml}: ClickHouse 26.3.12.3, порт 8124 только на loopback, база \texttt{mvb\_analysis}, пользователь \texttt{mvb}. Данные сохраняются в постоянном томе Docker. Пароль задают переменной среды \texttt{MVB\_CLICKHOUSE\_PASSWORD} через локальную конфигурацию и не включают в публикуемые материалы. Другой адрес, пользователь и база задаются переменными среды; для удалённого подключения с паролем требуется HTTPS.
\begin{verbatim}
docker compose -f infra/clickhouse.compose.yaml up -d
py -3 run_analysis.py --export-run UUID_СЕРИИ
  --output-dir output/reexport_001
\end{verbatim}
UUID заменяют значением из сводки первоначального запуска. В базе хранятся покадровые измерения, все обработанные кадры и протокол серии с калибровкой, параметрами, версиями и хешами. Исходное видео остаётся файлом. SQL вычисляет медианы, средние, разброс по траекториям и показатели кадров; графики строятся по повторно прочитанным данным базы. CSV, JSON, Markdown и графики являются производными экспортами. Повторный экспорт не требует видео, но аннотированные изображения без исходных кадров не восстанавливаются.

Каждый запуск получает новый UUID. Протокол завершения записывается после проверки числа измерений, кадров и диагностических записей. Частичные записи без него не доступны командой повторного экспорта; автоматического повторения вставок нет. При сбое базы серия не объявляется сохранённой. Для диагностики без сервера используют явный ключ \texttt{--offline}; этот режим не является основным порядком хранения. Старый режим \texttt{--single-ball} также требует \texttt{--offline}. Создание JSON калибровки по CSV остаётся вспомогательной операцией; применённый JSON сохраняется в протоколе серии в ClickHouse.

\section{Пакетная обработка нескольких шариков}
\begin{verbatim}
py -3 run_analysis.py --video data/videos/series.mp4
  --calibration calibration.json --measurement-mode reference
  --output-dir output/series_001 --frame-step 1
  --min-ball-diam-mm 1 --max-ball-diam-mm 10
  --tracking-distance-px 60 --tracking-max-missed 3
  --min-track-observations 2 --preview-frames 12
\end{verbatim}
Диапазон 1--10 мм и ограничения сопровождения являются стартовыми параметрами примера. Их устанавливают по ожидаемым размерам и перемещениям между \emph{обработанными} кадрами. Увеличение \texttt{frame-step} меняет интервал сопровождения и может приводить к потере идентичности. Каждый запуск получает новую папку вывода.

По умолчанию поиск ограничивается внутренней областью обнаруженной круглой ячейки. Если в кадре нет подходящего круглого референса, применяют \texttt{--no-cell-mask} после проверки области анализа. Это отключает ограничение ячейкой и требует контроля ложных окружностей в фоне. Ограничение числа обработанных кадров задаётся \texttt{--max-frames}. Сохранённая калибровка поддерживается многообъектным режимом; её нельзя сочетать с \texttt{--single-ball}. Интерактивный просмотр требует \texttt{--single-ball} и представляет отдельный прежний режим.

Алгоритм выполняет CLAHE и гауссово сглаживание, поиск окружностей Хафа, проверку краевой опоры, устранение дубликатов и уточнение окружности по краевым точкам исходного разрешения. Сопровождение использует глобальное сопоставление с прогнозом скорости и ограничением изменения размера. Идентификатор трека является результатом алгоритма, а не гарантированным номером физического шарика: после перекрытия или длительной потери один объект может получить новый номер.

\clearpage
\subsection{Устойчивое уточнение и отказы}
Для проверки контура предусмотрен дополнительный режим \texttt{--edge-fit-method radial}. Он использует радиальные градиенты и геометрическую подгонку с функцией потерь \texttt{soft\_l1}. Для тёмного шарика профиль \texttt{--detector-profile dark-ball} использует рост яркости наружу. Общий профиль допускает \texttt{--radial-polarity either}, \texttt{rising} и \texttt{falling}. Направление выбирают по исходным изображениям до проверочного прогона.

\begin{verbatim}
py -3 run_analysis.py --video data/videos/bearing.mp4
  --output-dir output/radial_001 --frame-step 1
  --measurement-mode pixel --detector-profile dark-ball
  --min-ball-diameter-px 140 --max-ball-diameter-px 250
  --edge-fit-method radial --edge-refine-band-px 25
  --radial-min-coverage 0.65 --radial-min-gradient 0.2
  --radial-max-axis-ratio 1.15 --no-cell-mask
  --annotated-video
\end{verbatim}
Числа задают пример для малой проекции в архивном 4К-видео; на другую оптику их не переносят без проверки. В новом режиме недостаточная угловая опора, большой остаток, эллиптичность и обрезанный контур вызывают отказ. Диагностика кандидатов, поступивших на уточнение, сохраняется в \texttt{quality.json} и в ClickHouse. Исходный кандидат вместо отклонённого измерения не выдаётся. Алгебраический режим \texttt{--edge-fit-method kasa} остаётся доступен для сравнения ; основной пакетный CLI выбирает \texttt{radial} по умолчанию. Выбор режима и порогов фиксируют в протоколе.

Сопровождение использует интервалы времени по меткам декодера. Число исправленных не возрастающих или недопустимых меток сохраняется в сводке. Видеофайл для просмотра имеет постоянную частоту; точное время измерения проверяют по CSV. Предсказанное положение без найденного контура не добавляется в таблицу измерений.

\section{Выходные материалы и контроль качества}
\begin{longtable}{p{48mm}p{82mm}}
\toprule Артефакт & Содержание и проверка\\\midrule\endhead
\path{measurements.csv} & Покадровые наблюдения: время, идентификатор, пиксельный и метрический диаметр; проверяют скачки и пригодность масштаба\\
\path{tracks.csv} & Число наблюдений, интервал сопровождения, медиана, среднее и разброс диаметра каждого трека\\
\path{frames.csv} & Число объектов и доступность масштаба по обработанным кадрам\\
\path{summary.json}, \path{summary.md} & Сводные результаты запуска\\
\path{quality.json} & Диагностика радиальной подгонки: принятые и отклонённые кандидаты, угловая опора и остаток; для старого режима список пуст\\
\path{run.json} & Параметры, сведения о среде и происхождении входа, SHA-256\\
\path{series_plots.png} & Временные зависимости и распределение медианных диаметров треков\\
Аннотированные кадры & Визуальная проверка попадания окружности на физическую границу\\\bottomrule
\end{longtable}
Гистограмма медиан треков уменьшает многократный учёт одного объекта в соседних кадрах. Однако повторное появление одного шарика с новым ID всё ещё может дать повторный учёт. Состав распределения проверяют по записям и траекториям.

Принимают серию к анализу после проверки масштаба, исходных кадров и устойчивости контуров. Причины исключения: насыщенный или размытый край, обмерзание окна, перекрытие, заметная деформация, ложная окружность на стенке, смена режима, непроверенный масштаб. Исключённые кадры и интервалы сохраняют в отдельном журнале с причиной; исходные данные и первоначальный экспорт не удаляют. Доля кадров с найденными объектами сама по себе не является полнотой детекции: для полноты требуется ручная или независимая разметка.

\section{Неопределённость и изменение уровня}
Для диаметра $d=NS$ в простой модели независимых вкладов
\[
u_{\mathrm{model}}(d)=\sqrt{(S u(N))^2+(N u(S))^2}.
\]
Экспортируемая модельная неопределённость не охватывает все аппаратные смещения. Для итогового измерения отдельно оценивают дисторсию, различие высот объекта и эталона, преломление, наклон, фокус, освещение, обмерзание и межсерийную повторяемость. Ковариации учитывают, если независимость вкладов не обоснована. Стандартную неопределённость не обозначают как гарантированный предел ошибки; при расширенной неопределённости указывают коэффициент охвата и обоснование.

Цель 0.01 мм из старого руководства не является достигнутой точностью. Её проверяют независимыми эталонами в рабочей среде и по всему используемому полю зрения. Субпиксельная подгонка контура также не доказывает аппаратную точность.

При изменении расстояния до объекта на $\Delta h$ простая перспективная модель даёт ориентир $\Delta S/S\approx\Delta h/D_0$. Старый пример $6.3/55\approx11.45\%$ показывает возможную значимость высоты, но не является измеренной поправкой данного стенда. В текущем сохранённом скалярном масштабе автоматической коррекции уровня нет. Работу с меняющимся уровнем допускают после определения поправки экспериментом либо при ограничении диапазона высот с оценкой остаточного вклада.

Для цилиндрической ячейки оценка испарившейся массы $m=\rho\pi D_{\mathrm{cell}}^2\Delta h/4$ требует независимо измеренных внутреннего диаметра, изменения уровня и плотности. Результат неприменим без проверки цилиндричности и условий измерения уровня. Управление насосом текущим пакетным анализом не реализовано; интерфейс, электрическая схема и обратная связь относятся к отдельному этапу.

\subsection{Таблица проверки подшипника}
Для завершённой серии одного подшипника выполняют:
\begin{verbatim}
python run_analysis.py --bearing-check-run UUID
  --bearing-reference-mm 3.98 --output-dir output/bearing_check
\end{verbatim}
Команду вводят одной строкой. UUID заменяют идентификатором серии из ClickHouse. Значение 3.98 мм предварительное; оно не задаёт неопределённость штангенциркуля. Первая половина записи формирует масштаб, вторая проверяет его перенос. Проверяют, что единственная траектория действительно относится к подшипнику.

Таблица содержит средний проверочный диаметр, условное смещение, RMSE, покадровый разброс и статистическую устойчивость среднего. Блочный bootstrap сохраняет временную сетку и пропуски, пересчитывает обе половины отдельно. Основной блок --- 15 кадров; результаты для 5 и 30 кадров сохранены в JSON. Не используют число кадров как число независимых повторных опытов.

Файлы \path{bearing_check.csv} и \path{bearing_uncertainty.json} формируются по измерениям из ClickHouse. Полная неопределённость сохраняется как NULL: для её оценки необходимы характеристики прибора, проверка масштаба независимым эталоном, дисторсии и рабочей оптической среды. Формулы и порядок оценки приведены в \path{docs/METROLOGY_BEARING_RU.md}.

\section{Проверки без камеры и последующая приёмка}
\begin{verbatim}
py -3 -m unittest discover -s tests -v
py -3 run_analysis.py --simulate-check
  --simulation-seed 907 --simulation-frames 24
  --simulation-output output/validation_new_903
\end{verbatim}
Повторный запуск использует новую папку результатов. На двух проверочных наборах seed 907 и 908 радиальная обработка обработала 480 синтетических кадров: 505 правильных обнаружений, 55 пропусков, 0 ложных обнаружений и 0 смен идентичности. Среди пропусков 48 относятся к деформированным объектам. Эти числа характеризуют выбранные размеченные сцены. Симуляция не подтверждает свойства реальной камеры, оптики, замерзание углеводородов или точность в жидком азоте. Подробные допущения зафиксированы в \path{docs/SIMULATION_METHOD_RU.md}.

Для аппаратной приёмки нужны: регистрация исходного видео без необъяснимых потерь; независимая эталонная проверка по полю зрения; повторение при перестановке и в рабочей среде; вручную размеченная выборка с пропусками, ложными обнаружениями и сменами ID; оценка неопределённости и сохранённый протокол. Целевые критерии и допустимые пределы согласуют до серии.

\section{Форма журнала серии}
\begin{enumerate}
\item ID серии, дата и время, оператор, цель и материал шариков.
\item Фотография монтажа; камера, объектив, кольца, рабочая геометрия и настройки регистрации.
\item Среда, уровень, измеренные размеры ячейки, состояние окна и эталон с неопределённостью.
\item Путь и SHA-256 исходного видео, файл калибровки, версия кода, точная команда запуска.
\item Папка результатов, визуальная проверка, исключённые интервалы и причины.
\item Итоговая выборка, показатели качества, модельная и расширенная неопределённости при наличии обоснования.
\end{enumerate}
\section{Архивные видео телефона и предварительная калибровка}
При отсутствии подтверждённого масштаба используют режим \texttt{pixel}; метрические столбцы сохраняются NULL. Для тёмного подшипникового шарика предусмотрен профиль \texttt{dark-ball}. Его нельзя автоматически переносить на светлый объект в кипящей жидкости. Пример для малой проекции на портретном 4К-видео:
\begin{verbatim}
py -3 run_analysis.py --video reference_phone.mp4
  --measurement-mode pixel --detector-profile dark-ball
  --min-ball-diameter-px 140 --max-ball-diameter-px 250
  --roi 0.25 0.30 0.75 0.65 --frame-step 1
  --no-cell-mask --annotated-video
  --output-dir output/phone_reference
\end{verbatim}
Диапазон и ROI являются примером конкретной архивной сцены. Размеры относятся к изображению после учёта поворота декодером. Контрольные JPG и MP4 дают форму результата; временные метки берут из таблицы измерений. В сводке проверяют число декодированных кадров и статус обрыва; ошибка NAL в исходнике не объявляется завершённой серией.

Если размер указан приблизительно, а неопределённость эталона неизвестна, допускается только предварительная калибровка:
\begin{verbatim}
py -3 run_analysis.py --calibrate-csv reference_one_track.csv
  --calibration-diameter-mm 3.98 --provisional-calibration
  --calibration-width-px 2160
  --calibration-output calibration_provisional.json
\end{verbatim}
Значение 3.98 мм приведено исполнителем для архивного подшипникового шарика, измеренного электронным штангенциркулем. Оно требует подтверждения протоколом. Полная неопределённость в этом режиме неизвестна; её нельзя заменять нулём. Проверка того же объекта, использованного для масштаба, не определяет абсолютную точность. Масштаб получают заново для каждого изменения фокуса, увеличения и геометрии; паспортное поле камеры телефона не заменяет такую калибровку.

\section{Источники}
При подготовке руководства использованы исходная инструкция по проведению эксперимента, документы на поставку оборудования и паспорт камеры \href{https://en.daheng-imaging.com/show-104-1901-1.html}{Daheng Imaging, MER2-231-41GM}. Программные процедуры описаны в \path{run_analysis.py} и \path{src/mvb/series_analysis.py}. Документация оборудования определяет параметры комплектации; условия конкретного опыта фиксируют в экспериментальном протоколе.
\end{document}
